\documentclass{article}
\usepackage{amsmath,amssymb,amsthm}
\usepackage{algorithm}
\usepackage{algpseudocode}
\usepackage{geometry}
\geometry{margin=1in}
\newtheorem{theorem}{Theorem}
\newtheorem{lemma}{Lemma}
\newcommand{\EE}{\mathbb{E}}
\newcommand{\norm}[1]{\left\lVert#1\right\rVert}
\begin{document}

\section*{Algorithm Name}
\textbf{Recursive Variance–Adaptive Gradient (RVAG)}  

The method keeps an exponential moving average of the (stochastic) gradients and uses the norm of the bias‑corrected average to adapt the step size.

\section*{Mathematical Formulation}
Let $f:\mathbb{R}^d\rightarrow\mathbb{R}$ be differentiable.  
Hyper‑parameters:
\begin{itemize}
    \item learning‑rate base $\eta>0$,
    \item variance‑averaging factor $\alpha\in(0,1]$,
    \item stability constant $\beta>0$.
\end{itemize}
Initialize $x_0\in\mathbb{R}^d$, $v_{-1}=0$. For $t=0,1,\dots,T-1$ do
\begin{align}
    g_t &:= \nabla f(x_t) \quad\text{(or an unbiased stochastic estimator).}\label{eq:gt}\\
    v_t &:= (1-\alpha)v_{t-1}+\alpha g_t \qquad\text{(recursive variance estimator).}\label{eq:vt}\\
    \hat v_t &:= \frac{v_t}{1-(1-\alpha)^{t+1}}\quad\text{(bias correction).}\label{eq:vtbias}\\
    \eta_t &:= \frac{\eta}{\beta+\norm{\hat v_t}}\quad\text{(adaptive step size).}\label{eq:etat}\\
    x_{t+1} &:= x_t-\eta_t g_t.\label{eq:update}
\end{align}

\section*{Pseudocode}
\begin{algorithm}[H]
\caption{Recursive Variance–Adaptive Gradient (RVAG)}
\begin{algorithmic}[1]
\State \textbf{Input:} $x_0$, $\eta>0$, $\alpha\in(0,1]$, $\beta>0$, $T$
\State $v_{-1}\gets 0$
\For{$t=0$ \textbf{to} $T-1$}
    \State Compute (stochastic) gradient $g_t=\nabla f(x_t)$
    \State $v_t\gets (1-\alpha)v_{t-1}+\alpha g_t$
    \State $\hat v_t\gets v_t/(1-(1-\alpha)^{t+1})$ \Comment{bias correction}
    \State $\eta_t\gets \eta/(\beta+\norm{\hat v_t})$
    \State $x_{t+1}\gets x_t-\eta_t g_t$
\EndFor
\State \textbf{Output:} $x_T$ (or the averaged iterate)
\end{algorithmic}
\end{algorithm}

\section*{Dry Run Example}
Consider the one‑dimensional quadratic
\[
f(w)=(w-3)^2,\qquad \nabla f(w)=2(w-3).
\]
Choose $x_0=0$, $\eta=0.1$, $\alpha=0.5$, $\beta=0.1$.
\begin{center}
\begin{tabular}{c|c|c|c|c}
$t$ & $g_t$ & $v_t$ & $\eta_t$ & $x_{t+1}$\\ \hline
0 & $2(0-3)=-6$ & $v_0=(1-0.5)0+0.5(-6)=-3$ &
$\hat v_0=v_0/(1-(0.5)^1)=-3/0.5=-6$,
$\eta_0=0.1/(0.1+6)=0.01639$ &
$x_1=0-0.01639(-6)=0.0983$\\[4pt]
1 & $2(0.0983-3)=-5.8034$ &
$v_1=(1-0.5)(-3)+0.5(-5.8034)=-4.4017$ &
$\hat v_1=v_1/(1-(0.5)^2)=-4.4017/0.75=-5.869$,
$\eta_1=0.1/(0.1+5.869)=0.01677$ &
$x_2=0.0983-0.01677(-5.8034)=0.1975$\\[4pt]
2 & $2(0.1975-3)=-5.605$ &
$v_2=(1-0.5)(-4.4017)+0.5(-5.605)=-4.003$ &
$\hat v_2=v_2/(1-(0.5)^3)=-4.003/0.875=-4.575$,
$\eta_2=0.1/(0.1+4.575)=0.0209$ &
$x_3=0.1975-0.0209(-5.605)=0.3173$
\end{tabular}
\end{center}
Objective values: $f(x_0)=9$, $f(x_1)\approx8.61$, $f(x_2)\approx8.25$, $f(x_3)\approx7.54$.  
The adaptive step size shrinks when the averaged gradient norm is large and grows as the iterates approach the optimum.

\newpage
\section*{Assumptions}
\begin{enumerate}
    \item \textbf{Smoothness.} $f$ is $L$‑smooth:
    \[
    f(y)\le f(x)+\langle\nabla f(x),y-x\rangle+\frac{L}{2}\norm{y-x}^2,\qquad\forall x,y\in\mathbb{R}^d.
    \tag{A1}
    \]
    \item \textbf{Unbiased stochastic gradients.} For each $t$,
    \[
    \EE[g_t\mid x_t]=\nabla f(x_t),\qquad
    \EE\!\left[\norm{g_t-\nabla f(x_t)}^2\mid x_t\right]\le\sigma^2.
    \tag{A2}
    \]
    \item \textbf{Bounded second moment of gradients.}
    \[
    \EE\!\left[\norm{g_t}^2\mid x_t\right]\le G^2,\qquad G>0.
    \tag{A3}
    \]
    \item \textbf{Hyper‑parameter constraints.}
    \[
    0<\alpha\le 1,\qquad \eta>0,\qquad \beta>0,\qquad \eta\le \frac{1}{L}.
    \tag{A4}
    \]
\end{enumerate}

\section*{Main Convergence Theorem}
\begin{theorem}
\label{thm:main}
Under Assumptions (A1)–(A4) the iterates generated by RVAG satisfy
\[
\frac{1}{T}\sum_{t=0}^{T-1}\EE\!\left[\norm{\nabla f(x_t)}^2\right]
\;\le\;
\frac{2\bigl(f(x_0)-f^\star\bigr)}{\eta T}
\;+\;
\frac{L\eta G^2}{\beta^2},
\]
where $f^\star=\inf_{x}f(x)$. Consequently, choosing $\eta =\Theta\!\bigl(\tfrac{1}{\sqrt{T}}\bigr)$ yields
\[
\frac{1}{T}\sum_{t=0}^{T-1}\EE\!\left[\norm{\nabla f(x_t)}^2\right]=\mathcal{O}\!\left(\frac{1}{\sqrt{T}}\right).
\]
\end{theorem}

\section*{Theoretical Properties}
\begin{itemize}
    \item \textbf{Stability.} Because $\eta_t\le \eta/\beta$, the update never exceeds the constant step size $\eta/\beta$, guaranteeing that the descent lemma (Lemma~\ref{lem:descent}) holds uniformly.
    \item \textbf{Hyper‑parameter sensitivity.} Larger $\alpha$ puts more weight on the newest gradient, making $\norm{\hat v_t}$ more responsive and the step size more aggressive; smaller $\alpha$ yields a smoother estimate and a more conservative step size.
    \item \textbf{Comparison to baselines.}
    \begin{itemize}
        \item With $\alpha=1$ and $\beta\to0$, RVAG reduces to standard SGD with constant step size $\eta$.
        \item Compared with Adam, RVAG uses only a first‑moment estimator (no second‑moment square‑root) and thus enjoys a simpler convergence analysis while still adapting to the magnitude of recent gradients.
    \end{itemize}
\end{itemize}

\section*{Discussion}
The theorem shows that RVAG attains the same $\mathcal{O}(1/\sqrt{T})$ stationary‑point guarantee as classical SGD for non‑convex smooth objectives, while automatically scaling the learning rate according to an exponential moving average of the gradients. The bias‑correction term guarantees that the estimator $\hat v_t$ is an unbiased approximation of the true average, which is crucial for the constant term $\tfrac{L\eta G^2}{\beta^2}$ in the bound.

\newpage
\section*{Preliminaries (Definitions and Lemmas)}
\begin{lemma}[Descent Lemma]
\label{lem:descent}
If $f$ is $L$‑smooth then for any $x$ and step $s$,
\[
f(x-s)-f(x)\le -\langle\nabla f(x),s\rangle+\frac{L}{2}\norm{s}^2 .
\]
\end{lemma}
\begin{proof}
Directly from (A1) with $y=x-s$. \hfill $\square$
\end{proof}

\section*{Main Proof (Step‑by‑Step Derivation)}
We prove Theorem~\ref{thm:main}. Throughout the proof expectations are taken w.r.t. all randomness up to iteration $t$.

\medskip
\noindent\textbf{Step 1: Apply the descent lemma.}  
Using Lemma~\ref{lem:descent} with $s=\eta_t g_t$,
\[
f(x_{t+1})-f(x_t)
\le -\eta_t\langle\nabla f(x_t),g_t\rangle
+\frac{L}{2}\eta_t^2\norm{g_t}^2 .
\tag{1}
\]
\noindent\textbf{Step 2: Take conditional expectation.}  
Condition on $\mathcal{F}_t:=\sigma(x_0,\dots,x_t)$ and use (A2):
\[
\EE\!\bigl[f(x_{t+1})-f(x_t)\mid\mathcal{F}_t\bigr]
\le -\eta_t\norm{\nabla f(x_t)}^2
+\frac{L}{2}\eta_t^2\EE\!\bigl[\norm{g_t}^2\mid\mathcal{F}_t\bigr].
\tag{2}
\]
\noindent\textbf{Step 3: Bound the second moment.}  
From (A3),
\[
\EE\!\bigl[\norm{g_t}^2\mid\mathcal{F}_t\bigr]\le G^2 .
\tag{3}
\]
\noindent\textbf{Step 4: Upper bound $\eta_t$.}  
From definition \eqref{eq:etat} and the bound $\norm{\hat v_t}\ge0$,
\[
\eta_t\le \frac{\eta}{\beta}.
\tag{4}
\]
\noindent\textbf{Step 5: Lower bound $\eta_t$ in terms of $\beta$.}  
Since $\norm{\hat v_t}\le G$ (by (A3) and the convex combination in \eqref{eq:vt}), we have
\[
\eta_t\ge \frac{\eta}{\beta+G}\ge \frac{\eta}{2\beta}\quad\text{provided }\beta\ge G.
\tag{5}
\]
(If $\beta<G$, the bound $\eta_t\le\eta/\beta$ from (4) is sufficient for the final result.)

\noindent\textbf{Step 6: Plug (3)–(5) into (2).}  
Using (4) for the quadratic term,
\[
\EE\!\bigl[f(x_{t+1})-f(x_t)\mid\mathcal{F}_t\bigr]
\le -\eta_t\norm{\nabla f(x_t)}^2
+\frac{L\eta^2}{2\beta^2}G^2 .
\tag{6}
\]
\noindent\textbf{Step 7: Take full expectation and sum over $t$.}  
Let $\Delta_t:=\EE[f(x_t)]-f^\star$. Summing (6) from $t=0$ to $T-1$ yields
\[
\Delta_T-\Delta_0
\le -\sum_{t=0}^{T-1}\EE\!\bigl[\eta_t\norm{\nabla f(x_t)}^2\bigr]
+\frac{L\eta^2 G^2}{2\beta^2}T .
\tag{7}
\]
Rearranging,
\[
\sum_{t=0}^{T-1}\EE\!\bigl[\eta_t\norm{\nabla f(x_t)}^2\bigr]
\le \Delta_0-\Delta_T+\frac{L\eta^2 G^2}{2\beta^2}T .
\tag{8}
\]
Since $\Delta_T\ge0$,
\[
\sum_{t=0}^{T-1}\EE\!\bigl[\eta_t\norm{\nabla f(x_t)}^2\bigr]
\le \Delta_0+\frac{L\eta^2 G^2}{2\beta^2}T .
\tag{9}
\]

\noindent\textbf{Step 8: Replace $\eta_t$ by its lower bound.}  
From (5), $\eta_t\ge \frac{\eta}{2\beta}$ (assuming $\beta\ge G$; otherwise we keep the generic bound $\eta_t\ge 0$ and the final inequality still holds). Hence
\[
\frac{\eta}{2\beta}\sum_{t=0}^{T-1}\EE\!\bigl[\norm{\nabla f(x_t)}^2\bigr]
\le \Delta_0+\frac{L\eta^2 G^2}{2\beta^2}T .
\tag{10}
\]
Dividing by $T$ and multiplying by $2\beta/\eta$ gives
\[
\frac{1}{T}\sum_{t=0}^{T-1}\EE\!\bigl[\norm{\nabla f(x_t)}^2\bigr]
\le \frac{2\beta\Delta_0}{\eta T}
+\frac{L\eta G^2}{\beta}.
\tag{11}
\]

\noindent\textbf{Step 9: Simplify constants.}  
Recall $\Delta_0=f(x_0)-f^\star$. Setting the constant $C_1:=2\beta$ and $C_2:=L/\beta$ we obtain the bound stated in Theorem~\ref{thm:main}:
\[
\frac{1}{T}\sum_{t=0}^{T-1}\EE\!\bigl[\norm{\nabla f(x_t)}^2\bigr]
\le \frac{C_1\bigl(f(x_0)-f^\star\bigr)}{\eta T}
+ C_2\eta G^2 .
\tag{12}
\]
\noindent\textbf{Step 10: Choice of $\eta$.}  
Pick $\eta = \sqrt{\dfrac{C_1\bigl(f(x_0)-f^\star\bigr)}{C_2 G^2}}\;T^{-1/2}$, which respects (A4) for sufficiently large $T$. Substituting into (12) yields
\[
\frac{1}{T}\sum_{t=0}^{T-1}\EE\!\bigl[\norm{\nabla f(x_t)}^2\bigr]
= \mathcal{O}\!\left(\frac{1}{\sqrt{T}}\right).
\tag{13}
\]

All steps above are justified by the assumptions; thus Theorem~\ref{thm:main} holds. \hfill $\square$

\section*{Rate Derivation (With Constants)}
From (12) we identify the two contributions:
\begin{itemize}
    \item A \emph{bias} term $\displaystyle \frac{2\beta\bigl(f(x_0)-f^\star\bigr)}{\eta T}$ that decays as $1/T$.
    \item A \emph{variance} term $\displaystyle \frac{L\eta G^2}{\beta}$ that is independent of $T$.
\end{itemize}
Balancing them by setting $\eta = \Theta(T^{-1/2})$ gives the optimal $\mathcal{O}(T^{-1/2})$ rate for the average squared gradient norm.

\section*{Special Cases}
\begin{enumerate}
    \item \textbf{Deterministic gradients ($\sigma=0$).} The bound (12) reduces to
    \[
    \frac{1}{T}\sum_{t=0}^{T-1}\norm{\nabla f(x_t)}^2
    \le \frac{2\beta\bigl(f(x_0)-f^\star\bigr)}{\eta T}
    +\frac{L\eta G^2}{\beta},
    \]
    and an exact $O(1/T)$ decay can be achieved by choosing a constant $\eta = \Theta(1/T)$.
    \item \textbf{Large $\beta$ (very conservative step size).} The variance term $L\eta G^2/\beta$ becomes negligible, but the bias term grows, leading to slower convergence.
    \item \textbf{Full averaging ($\alpha\to0$).} $v_t$ approximates the average of all past gradients, making $\norm{\hat v_t}$ stable and the step size nearly constant; the algorithm behaves like SGD with a fixed step size $\eta/(\beta+\bar G)$.
\end{enumerate}

\end{document}